\begingroup
\small\setstretch{1.12}\setlength{\tabcolsep}{3pt}
\setlength{\LTleft}{\fill}\setlength{\LTright}{\fill}
\Needspace{17\baselineskip}
\begin{longtable}{@{}>{\raggedright\arraybackslash}p{0.25\linewidth}>{\raggedright\arraybackslash}p{0.69\linewidth}@{}}
\caption{本研究本機設備與固定年份執行設定}\label{tab:method-hardware}\\
\toprule
項目 & 規格／設定 \\
\midrule\endfirsthead
\multicolumn{2}{c}{表~\thetable（續）}\\
\toprule
項目 & 規格／設定 \\
\midrule\endhead
\midrule\multicolumn{2}{r}{續下頁}\\\endfoot
\bottomrule\endlastfoot
CPU & Intel Core i7-10870H，2.20 GHz；8 核心、16 邏輯處理器 \\
記憶體 & 32 GB（安裝容量） \\
作業系統 & Windows 11 家用版，64 位元；版本 10.0.26200 \\
獨立顯示卡 & NVIDIA GeForce GTX 1650 Ti；4096 MiB 顯示記憶體 \\
內建顯示晶片 & Intel UHD Graphics \\
本機 NVIDIA 驅動 & 592.82（本次查核值） \\
固定年份 XGBoost & CPU；hist 樹建構，未設定 CUDA device \\
FS 執行 & scikit-learn 決策樹／MI 與 LightGBM CPU 設定；最多四執行緒 \\
查核日期 & 2026 年 10 月 27 日 \\
\end{longtable}
\noindent{\footnotesize (資料來源：依使用者指定，採本機系統查核與實驗程式設定整理)\par 設備具有 GPU 不代表所有方法都使用 GPU；本次查得的作業系統及驅動版本不回推為每次歷史執行的原始版本。}
\endgroup
